% MAS1004 Assignment 1: long report.
%
% You fill this in with your agent, in your repository. AGENTS.md tells the
% agent how. Red text in the PDF is what still has to be replaced.
%
% - Pictures go in figures/. Copy them from results/, which is not in git.
% - Outputs go in outputs/. run.py saves what it prints in
%   results/<tag>_output.txt, and PLAN.md saves the other outputs in results/
%   with `| tee results/...`. Copy them here. \outputfile{outputs/check.txt}
%   prints a file exactly as it is. Every number in this report has to come
%   from such an output.
%
% To make the PDF: compress this folder into a zip, and on Overleaf choose
% New Project, Upload Project. Set Menu, Compiler to XeLaTeX, and Recompile.

\documentclass[11pt]{article}
\usepackage{mas1004}

\reporttitle{Assignment 1: Long report}
\studentname{Ten Aleksandra}
\studentid{20261332}

\begin{document}
\makeheader

% ---------------------------------------------------------------------------
\section*{Problem 1. Categories and images}

\subsection*{Why these categories}
I am a big fan of Sanrio and generally love cute things, so I wanted to make
the project a little cuter. I also thought the characters were visually
distinct enough that training would go smoothly.

\subsection*{What I expected the model to use}
Before training, I expected colour schemes and distinctive shapes to be the
most obvious differences between the characters.

\subsection*{Images per category}
The downloaded set had Cinnamoroll: 150, Hello Kitty: 107, Kuromi: 150, and
Pompompurin: 150 images.

% ---------------------------------------------------------------------------
\section*{Problem 2. Preparing an image}

\sidebyside{figures/prepare_original.jpg}{Original Pompompurin image}
           {figures/prepare_result.png}{The centre crop made by \texttt{prepare\_image}}

The centre crop removes part of the small figure and surrounding context at
the lower left. The main Pompompurin figure remains in the crop.

% ---------------------------------------------------------------------------
\section*{Problem 3. First training runs}

\begin{center}\footnotesize
\begin{tabular}{llllrrrr}
\toprule
tag & start & trained & epochs & lr & trainable parameters & train acc. & test acc. \\
\midrule
run & ImageNet & all layers & 10 & 0.0001 & 11,178,564 & 100.0\% & 97.3\% \\
scratch & random & all layers & 10 & 0.0001 & 11,178,564 & 98.7\% & 81.1\% \\
frozen & ImageNet & last layer & 10 & 0.001 & 2,052 & 96.4\% & 92.8\% \\
short & ImageNet & all layers & 3 & 0.0001 & 11,178,564 & 100.0\% & 98.2\% \\
clean & ImageNet & all layers & 10 & 0.0001 & 11,178,564 & 100.0\% & 98.9\% \\
\bottomrule
\end{tabular}
\end{center}

The run loss fell quickly and then flattened; training loss continued to fall
while test loss fluctuated near its later values. Training accuracy was
100.0\% and test accuracy was 97.3\%, a small train-test gap. In the run
confusion matrix, the largest confusion was Cinnamoroll predicted as Hello
Kitty. This was plausible given my expectation that the model would use
visual appearance such as colour and shape.

Starting from ImageNet weights gave a higher test accuracy than starting from
random weights: 97.3\% compared with 81.1\%. Training only the last layer
reached 92.8\%, below the all-layer run's 97.3\%, but still relatively close.
The ImageNet model already had useful visual features, while the model trained
from scratch had to learn them from this smaller assignment dataset.

\begin{center}
\picturefile[0.48\linewidth]{figures/run_curves.png}\hfill
\picturefile[0.48\linewidth]{figures/run_confusion.png}
\end{center}
\begin{center}
\picturefile[0.48\linewidth]{figures/scratch_curves.png}\hfill
\picturefile[0.48\linewidth]{figures/scratch_confusion.png}
\end{center}
\begin{center}
\picturefile[0.48\linewidth]{figures/frozen_curves.png}\hfill
\picturefile[0.48\linewidth]{figures/frozen_confusion.png}
\end{center}
\begin{center}
\picturefile[0.48\linewidth]{figures/short_curves.png}\hfill
\picturefile[0.48\linewidth]{figures/short_confusion.png}
\end{center}
\begin{center}
\picturefile[0.48\linewidth]{figures/clean_curves.png}\hfill
\picturefile[0.48\linewidth]{figures/clean_confusion.png}
\end{center}

\subsection*{Output of the runs}
\outputfile{outputs/run_output.txt}
\outputfile{outputs/scratch_output.txt}
\outputfile{outputs/frozen_output.txt}
\outputfile{outputs/short_output.txt}

% ---------------------------------------------------------------------------
\section*{Problem 4. Cleaning}

\subsection*{My rule}
My rule was: ``remove an image if more than one Sanrio character is clearly
visible, or it is fan art or a different art style, or the character is only
a small secondary element (sticker or logo on merchandise), not the main
subject; keep it if exactly one character is the clear main subject''.

\subsection*{Three kinds of junk}
Examples of the kinds I removed:

\begin{center}
\picturefile[0.31\linewidth]{figures/junk_multiple_characters.jpg}\hfill
\picturefile[0.31\linewidth]{figures/junk_fan_art.jpg}\hfill
\picturefile[0.31\linewidth]{figures/junk_secondary.jpg}
\end{center}
\begin{center}
Multiple characters \hfill Fan art or a different style \hfill
Character secondary to merchandise
\end{center}

\subsection*{What I removed}
\outputfile{outputs/clean_output.txt}
\outputfile{outputs/clean_count.txt}

The removal log records 22 images removed as near copies. I stopped before a
second pass over Kuromi and Pompompurin, so some near copies are still in the
data. I also kept the near-copy sheet for reference:
\picturefile[0.8\linewidth]{figures/copies.png}

Among the first 20 suspects in each class, I removed 12 Cinnamoroll, 8 Hello
Kitty, 7 Kuromi, and 2 Pompompurin images. The suspect list revealed junk I
had missed, and my manual review also found junk that was not listed. The
suspect counts were matched from the saved suspect ranking and removal log.

\subsection*{Before and after cleaning}
\outputfile{outputs/compare.txt}

On the downloaded test set, accuracy rose from 97.3\% before cleaning to
98.9\% after cleaning. On the unchanged new images it rose from 82.1\% to
89.3\%. Neither accuracy went down; removing images outside my category rule
may have made the training data more consistent.

% ---------------------------------------------------------------------------
\section*{Problem 5. Images from a new source}

I took some of the new images myself and took screenshots from unpopular
videos. They were not collected by the web search used for the downloads.

\outputfile{outputs/overlap.txt}

The cleaned model's downloaded-test accuracy was 98.9\%; its accuracy on the
new images was 89.3\%. The new images had different lighting and less clear
backgrounds, which I think contributed to the mistakes.

The saved check output lists three confident mistakes on the new images, so I
include those three rather than inventing more. The downloaded-set mistake is
shown separately in the saved worst-example plot: a Hello Kitty image was
predicted as Cinnamoroll.

\picturefile[0.72\linewidth]{figures/clean_worst.png}

\begin{center}
\picturefile[0.31\linewidth]{figures/20261015_181714.jpg}\hfill
\picturefile[0.31\linewidth]{figures/20261015_183443.jpg}\hfill
\picturefile[0.31\linewidth]{figures/20261015_183901.jpg}
\end{center}
\begin{center}
Hello Kitty was called Pompompurin; Cinnamoroll was called Kuromi;
Kuromi was called Cinnamoroll. Differences in lighting and unclear
backgrounds may have contributed.
\end{center}

% ---------------------------------------------------------------------------
\section*{Problem 6. The demo}

Address: \url{https://zunogi.github.io/sanrio-classifier/}

The demo still assigns an input outside the categories to one of its four
classes.

% ---------------------------------------------------------------------------
\section*{Problem 7. What the model looks at}

The guesses came from quick manual tests on the demo page, including changing
images to grayscale and covering parts. I was also suspicious that training
had gone too smoothly.

\subsection*{Exact request to the agent}
\begin{Verbatim}[fontsize=\footnotesize,breaklines]
Part 1. Do not modify anything in data/new_images (the 28 originals).
Create copies under data/changed/, keeping the same class sub-folders
(data/changed/<experiment>/<class>/<same file name>):
1. grayscale: convert each image to grayscale, saved as 3-channel RGB.
2. hue180: shift the hue of each image by 180 degrees (same shape, different colours).
3. top_third_grey: only for the images in data/new_images/pompompurin, cover the top third
   of the image with a flat grey (128,128,128) rectangle.
4. bottom_third_grey: same, but cover the bottom third instead.
5. tile2x2: for every image in data/new_images, make one image that is the same picture
   repeated in a 2x2 grid (same total size as the original).
Do not change any file in src/.

Part 2. Write a new script src/experiments.py. It loads the trained model from
results/clean_model.pt (use train.load), and for data/new_images and for each folder in
data/changed it uses data.load_folder to read the images and predicts with softmax.
For each folder print and save to results/experiments.csv:
accuracy, how many images were predicted as each class, and the mean probability given
to the TRUE class. Also save results/experiments_per_image.csv with one row per image:
experiment, file, true class, predicted class, confidence, probability of true class.
Do not change any existing function. Show me the printed table.
\end{Verbatim}

\subsection*{Guess 1: dark or grayscale images are called Kuromi}
I tested grayscale images because quick demo tests had suggested that losing
colour could change the answer. Grayscale tests whether the model relies on
colour; I did not separately darken images, so the dark-image part of the
guess remains untested. Accuracy was 0.4286 on grayscale images, and the model
predicted Kuromi for 25 images.

\sidebyside{figures/guess1_original.jpg}{Original answer: Cinnamoroll}
           {figures/guess1_grayscale.jpg}{Grayscale answer: Kuromi}

\subsection*{Guess 2: colour matters}
I thought yellow might matter for Pompompurin and light blue for Cinnamoroll.
The hue shift tests whether changing colour changes predictions, but it
changes the colours across the whole image and does not isolate those exact
colour cues. Accuracy was 0.7143 after the hue shift.

\sidebyside{figures/guess2_original.jpg}{Original answer: Cinnamoroll}
           {figures/guess2_hue180.jpg}{Hue-shifted answer: Kuromi}

\subsection*{Guess 3: Pompompurin is recognised by its hat}
Covering the top third was intended to test whether removing the region with
the hat changes Pompompurin predictions. The top-third experiment's accuracy
was 0.8333: 5 of the 6 Pompompurin images remained classified as Pompompurin.
One example changed from Pompompurin to Kuromi. This guess was not confirmed:
the mask covered more than just the hat, and the test does not isolate the hat
as the cause.

\sidebyside{figures/guess3_original.jpg}{Original answer: Pompompurin}
           {figures/guess3_top_third_grey.jpg}{Top third grey answer: Kuromi}

\subsection*{Guess 4: repeated characters confuse the model}
Repeating each image in a 2x2 grid tests the effect of seeing several smaller
copies of the same character. Accuracy was 0.6429 after tiling. The lower
accuracy is consistent with the guess, although tiling also makes each
character smaller.

\sidebyside{figures/guess4_original.jpg}{Original answer: Cinnamoroll}
           {figures/guess4_tile2x2.jpg}{Tiled answer: Pompompurin}

\subsection*{What I now think the model uses}
Colour and tone seem to be the main factors to me because grayscale images
had the worst results. I am not certain that colour alone explains the
predictions: the colour changes affected whole images, and the top-third and
tiling tests changed more than one visual property.

\subsection*{Experiment summaries and per-image results}
\outputfile{outputs/experiments.txt}
\outputfile{outputs/experiments.csv}
\outputfile{outputs/experiments_per_image.csv}

% ---------------------------------------------------------------------------
\section*{Appendix. Output of check.py}
\outputfile{outputs/check.txt}

\end{document}
